% !TeX root = ../抽象代数笔记.tex
% 由主文档通过 \input 载入；请编译主文档。
\setcounter{chapter}{7} % 保留原书章号。
\chapter{直积与有限阿贝尔群}\label{ch:abelian}

\begin{description}
  \item[学习目标] 理解群的分解，掌握有限阿贝尔群分类及其同构判别。
  \item[先修知识] 第\ref{ch:quotient-groups}章；有限群相关例子结合第\ref{ch:sylow}章。
\end{description}

\section{外直积与内直积}
定义直积（direct product）$G\times H$ 的逐分量运算。对群 $G$ 的两个正规子群 $H,K$，在 $G=HK$ 且 $H\cap K=\{e\}$ 时证明 $G\cong H\times K$；先证明两子群的元素彼此交换，再验证乘法映射为同构。推广到有限个因子的情形。

\section{循环群的直积分解}
证明 $\mathbb Z/m\mathbb Z\times\mathbb Z/n\mathbb Z$ 为循环群当且仅当 $\gcd(m,n)=1$。用素数幂因子分解有限循环群，比较 $\mathbb Z/12\mathbb Z$ 与 $\mathbb Z/6\mathbb Z\times\mathbb Z/2\mathbb Z$ 的元素阶。将这类分解与第\ref{ch:ideals}章的环版本中国剩余定理衔接。

\section{有限阿贝尔群结构定理}
陈述并证明有限阿贝尔群是有限个素数幂阶循环群的直积，且这些因子除排列顺序外唯一；再改写为非平凡循环因子的阶逐项整除的形式。分别引入初等因子（elementary divisors）和不变因子（invariant factors）。证明路线依次为素数分量分解、阿贝尔 $p$ 群的循环直因子分解、由各 $p^k$ 阶挠子群大小恢复因子重数，不依赖后面的模论。

\section{半直积初步（选学）}
定义自同构群 $\operatorname{Aut}(N)$。给定群同态 $\varphi:H\to\operatorname{Aut}(N)$，在 $N\times H$ 上用 $(n,h)(n',h')=(n\varphi(h)(n'),hh')$ 定义半直积（semidirect product）$N\rtimes_\varphi H$，验证群公理。用二面体群说明正规因子与补子群的元素不必彼此交换，直积只对应作用平凡的特例。

\section{例题、练习与自检}
列出所有 $72$ 阶阿贝尔群的同构类型，转换初等因子与不变因子表示，并用元素阶或挠子群大小区分同阶群。比较 $S_3$ 的半直积描述与 $\mathbb Z/6\mathbb Z$ 的直积描述；强调有限阿贝尔群分类不适用于一般有限群。

